%=======================================================================================
% DATA CHAPTER -- rewritten 2026-08-12, following an approved chapter-plan/fact-inventory pass.
%
% Base content: the pasted draft supplied for that planning pass (not identical to either
% data_chapter_draft.tex or data_chapter_merged_2026-08-10.tex already on disk -- this file
% supersedes both for the 2026-08-12 structure; the two earlier files are left untouched in case
% any of their inline "Confirmed"/"NEEDS CHECKING" annotations are still wanted for reference).
%
% Every number below was re-verified this pass directly against data_processing/clean_master_data.py
% (re-executed live against the raw GeoPackage, not just read), the processed Parquet tables, and
% documentation/methodlogy_env_setpick.md (env. variables, verified 2026-08-10). Colour-coded
% flags follow the same legend already established in 1_Chapter_Intro/11th_aug_introdraft.tex:
% \fc = fact-check needed (red), \ct = citation missing (orange), \app = move to appendix (blue),
% \sr = small reword (teal). No model results, performance numbers, or causal language appear in
% this chapter -- see Methodology/Results for those.
%=======================================================================================
\newcommand{\fc}[1]{\textcolor{red}{[FactC: #1]}}
\newcommand{\ct}[1]{\textcolor{orange}{[CITE: #1]}}
\newcommand{\app}[1]{\textcolor{blue}{[APP: #1]}}
\newcommand{\sr}[1]{\textcolor{teal}{[Reword: #1]}}
%=======================================================================================

\chapter{Data}
\label{ch:data}

This chapter describes the data on which every result in this dissertation is built: where they
came from, how the analytical cohorts were constructed, what the response variable and the
original stand-management fields mean, and what external terrain, wind, climate and soil data
were attached to each plot. It states what exists and how it was assembled, not how it was
selected, split or modelled for a specific research question -- those decisions are described in
Chapter~\ref{ch:methodology}.

\section{Study area and source data}
\label{sec:data-source}

The plot data used throughout this dissertation were supplied by Forest Research as a single
GeoPackage layer (\texttt{LiDAR\_Years}), covering one study area, Aberfoyle forest in the
Trossachs, Scotland \ct{Aberfoyle forest location/general description -- a short geographic
citation, not yet in the citation list}. The \texttt{AOI} field is constant across every one of
the 795,645 raw rows (\texttt{"Aberfoyle"}), confirming a single, unambiguous study area rather
than a merged multi-site extract. Coordinates are recorded in the British National Grid
(EPSG:27700); the plots used in the final analytical cohorts span roughly 16.8\,km east--west and
12.0\,km north--south.

The raw layer holds 795,645 plot-year records across 230,359 unique plot identifiers and 38
columns, spanning six survey years: 2002 (37,888 records), 2006 (22,538), 2008 (114,975), 2012
(173,550), 2021 (218,359) and 2023 (228,335). Coverage grows substantially between 2006 and 2008,
consistent with an expansion of the sampling grid rather than plot loss.

Each record represents one plot in one survey year. A plot is the individual Area-Based-Analysis
sampling grid cell over which per-survey statistics (height percentiles, canopy cover, and the
other fields listed below) were computed from Airborne Laser Scanning (ALS) data -- an
aircraft-mounted system, not a drone survey. The nominal grid cell is 20\,m by 20\,m; the
\texttt{area} field's median is exactly 400\,m\textsuperscript{2} (71.1\% of all raw rows,
565,482 of 795,645, are exactly 400\,m\textsuperscript{2}), with a lower tail down to
250.0\,m\textsuperscript{2} where a grid cell is clipped by a compartment or sub-compartment
boundary. \fc{The area field also has an upper tail, up to 1{,}585.9\,m\textsuperscript{2} (almost
four times the nominal cell), which the boundary-clipping explanation does not account for --
clipping can only shrink a cell, not enlarge it. This needs checking against Forest Research's own
plot-layout documentation before the chapter states a cause; until then, report only the median
and the clipped lower tail with confidence, and flag the upper tail as unexplained.} Exact
acquisition metadata for the ALS survey itself (point density, sensor, flight dates) was not found
in any file inspected for this chapter; \ct{Forest Research ALS/LiDAR acquisition specification,
if a technical report exists} -- otherwise this remains a stated limitation of the source data's
documentation, not an assumption.

\app{Full raw column inventory (all 38 fields in \texttt{LiDAR\_Years}), including fields never
used downstream (e.g.\ \texttt{flyr}, constant at zero for every row; \texttt{block}, a
99.999\%-identical duplicate of \texttt{blk}).}

\section{Spatial and temporal hierarchy}
\label{sec:data-hierarchy}

The plot data sit inside a three-level Forest Research management hierarchy, and the naming of
that hierarchy in the raw layer is not as simple as the field names suggest. \texttt{blk} (block)
is the coarsest unit: 18 distinct blocks, and a forestry compartment (\texttt{cpmt}, 531 distinct
values in the raw layer) sits inside exactly one block in all but 2 of the 531 cases
(\fc{the two compartments spanning more than one block have not been individually checked -- worth
a footnote confirming whether this reflects a genuine boundary redefinition or a data-entry
artefact}). Below compartment, the raw layer's \texttt{scpt} field is \emph{not} a unique
sub-compartment identifier on its own: it takes only 26 distinct values, and inspection shows they
are simply the letters A--Z, reused across many different compartments (every one of the 26
letters occurs in more than one compartment). The genuine sub-compartment unit is therefore the
combination \texttt{(cpmt, scpt)}, which yields 3,385 distinct sub-compartments in the raw layer.
A plot is the finest unit still: the individual ALS sampling grid cell that survey statistics were
computed over, with many plots sitting inside one sub-compartment.

The six survey years are unevenly spaced: 2002 to 2006 (4 years), 2006 to 2008 (2 years), 2008 to
2012 (4 years), 2012 to 2021 (9 years), and 2021 to 2023 (2 years). This irregularity, together
with the coverage expansion after 2006, motivates treating the six-survey and four-survey cohorts
(\S\ref{sec:data-cohorts}) as different sampling problems rather than one dataset with some rows
missing.

% FIGURE PLACEHOLDER (study area and coverage)
% Caption: "Aberfoyle study area (British National Grid, EPSG:27700) with compartment boundaries
% and plot coverage for the four-survey cohort (left) and the six-survey cohort (right), shown at
% matched map scale. The six-survey cohort's plots are a complete subset of the four-survey
% cohort's plots (Table [cohort comparison]); the map shows which part of the forest that
% requirement retains."
% Content: two-panel map, national/regional inset in one corner, compartment outlines from
% data/interim/compartment_boundaries.parquet, plot points or a density surface (not raw scatter,
% to avoid overplotting at this record count) from each cohort's own master parquet.
\begin{figure}[!htbp]
\centering
\fbox{\parbox{0.9\textwidth}{\centering\vspace{4cm}FIGURE PLACEHOLDER --- study area and cohort
coverage, see comment above source for exact content\vspace{4cm}}}
\caption{Study area, spatial hierarchy and cohort coverage.}
\label{fig:data-study-area}
\end{figure}

\section{Cohort construction and coverage}
\label{sec:data-cohorts}

Two analytical cohorts were built from the same raw layer, differing only in which survey years
are required to be present for every retained plot: a \textbf{four-survey} cohort (2008, 2012,
2021, 2023) and a \textbf{six-survey} cohort (2002, 2006, 2008, 2012, 2021, 2023). Both are
produced by the same six-stage cleaning funnel, re-run directly against the current raw GeoPackage
for this chapter and matching the pipeline's own printed output exactly:

\begin{enumerate}
\setlength\itemsep{0pt}
\item Restrict to the cohort's chosen survey years.
\item Keep only plots present in \emph{every} one of those years (a whole plot is dropped if any
  one of its required survey rows is missing -- balanced trajectories only, never partial ones).
\item Keep only Sitka spruce, using the Forest Research inventory species code
  \texttt{spis == "SS"}. Two other species-related fields exist in the raw layer
  (\texttt{Species}, \texttt{SPIS\_RF}) but were not used for filtering: they carry only 9 coarse
  categories each (versus \texttt{spis}'s 41), are derived from a different source pipeline, and
  disagree with \texttt{spis} on a substantial share of rows.
\item Keep only a valid planting year (\texttt{plyr} $\geq$ 1850 and strictly before the survey
  year).
\item Keep only a plausible age (\texttt{Age}, survey year minus planting year, between 0 and 200
  years).
\item Keep only a plausible height (the target column, \texttt{elev\_percentile\_95th}, greater
  than 0\,m and at most 60\,m).
\end{enumerate}

Stages 4--6 are wide sanity bounds, not domain-derived thresholds, and on the current raw data
they remove very little: stage 4 removes rows only from the four-survey cohort, and stages 5--6
remove none at all from either cohort -- almost all attrition happens at stages 2--3 (requiring
complete coverage, then requiring Sitka spruce).

\begin{table}[!htbp]
\centering
\scriptsize
\setlength{\tabcolsep}{3pt}
\renewcommand{\arraystretch}{0.95}
\begin{tabular}{p{4.6cm}p{2.6cm}p{2.6cm}p{2.6cm}p{2.6cm}}
\hline
Stage & Four-survey rows & Four-survey plots & Six-survey rows & Six-survey plots \\
\hline
0. Raw dataset & 795,645 & 230,359 & 795,645 & 230,359 \\
1. Cohort years & 735,219 & 230,345 & 795,645 & 230,359 \\
2. Complete survey coverage & 451,128 & 112,782 & 112,686 & 18,781 \\
3. Sitka spruce & 289,580 & 72,395 & 83,382 & 13,897 \\
4. Valid planting year & 287,064 & 71,766 & 83,382 & 13,897 \\
5. Plausible age & 287,064 & 71,766 & 83,382 & 13,897 \\
6. Plausible height & 287,064 & 71,766 & 83,382 & 13,897 \\
\hline
\end{tabular}
\caption{Cleaning funnel, re-executed directly against the current raw GeoPackage
(\texttt{data\_processing/clean\_master\_data.py::clean\_cohort}, verified 2026-08-12).}
\label{tab:data-funnel}
\end{table}

\begin{table}[!htbp]
\centering
\scriptsize
\setlength{\tabcolsep}{2.5pt}
\renewcommand{\arraystretch}{0.9}
\begin{tabular}{p{4.3cm}p{3.2cm}p{3.2cm}}
\hline
Characteristic & Four-survey cohort & Six-survey cohort \\
\hline
Final plots & 71,766 & 13,897 \\
Final plot-year records & 287,064 & 83,382 \\
Survey years & 2008, 2012, 2021, 2023 & 2002, 2006, 2008, 2012, 2021, 2023 \\
Observations per plot & 4 & 6 \\
Compartments (\texttt{cpmt}) covered & 296 & 48 \\
Role & Main analysis & Sensitivity check \\
\hline
\end{tabular}
\caption{Cohort comparison. \fc{Six-survey compartment count verified directly against
\texttt{clean\_master\_6survey.parquet} and \texttt{growth\_curve\_table.parquet} (both give 48),
and matches \texttt{models/common/splits.py}'s own comment and the Results chapter draft. The
current Methodology draft (\S\ref{sec:rq2}) states 47 -- this appears to be inherited from an
older, pre-rebuild pipeline run recorded in \texttt{experiment\_log.md} whose plot/compartment
counts (13,769 plots / 47 compartments) do not match the current pipeline at all. Recommend
updating Methodology to 48 rather than re-deriving a new number here.}}
\label{tab:data-cohort-comparison}
\end{table}

The six-survey cohort is a \textbf{complete subset} of the four-survey cohort: every one of its
13,897 plots also appears in the four-survey cohort (0 six-survey-only plots), so the two cohorts
are nested rather than partially overlapping. This nesting is what makes the six-survey cohort a
meaningful sensitivity check on the same underlying plots, rather than an independent replication
-- how that check is used for a given research question is described in
Chapter~\ref{ch:methodology}.

% FIGURE PLACEHOLDER (funnel + timeline)
% Caption: "Left: cleaning funnel from the raw 795,645-record layer to the two final cohorts,
% showing that most attrition happens at the complete-coverage and species-filter stages, not the
% plausibility checks. Right: the six survey years plotted on a real timeline, showing the
% irregular gaps (4, 2, 4, 9, 2 years) and which years each cohort uses."
% Content: left panel, a stepped/Sankey-style bar from Table funnel's numbers, two parallel paths
% from one raw source; right panel, a horizontal timeline with two rows (four-survey / six-survey)
% and tick marks at the six real years, gap lengths labelled.
\begin{figure}[!htbp]
\centering
\fbox{\parbox{0.9\textwidth}{\centering\vspace{4cm}FIGURE PLACEHOLDER --- cleaning funnel and
survey timeline, see comment above source for exact content\vspace{4cm}}}
\caption{Cohort-construction funnel and repeated-observation timeline.}
\label{fig:data-funnel-timeline}
\end{figure}

\section{Response variable and original stand fields}
\label{sec:data-response}

The response variable used throughout this dissertation is \texttt{elev\_percentile\_95th}: the
95th percentile of ALS return heights above ground within a plot, for a given survey year, kept
in its raw form with no correction applied. A related field, \texttt{Top\_Height95}, is exactly
\texttt{elev\_percentile\_95th}~$\times$~1.1 (confirmed exactly on the raw layer, standard
deviation of the ratio $2\times10^{-17}$) and is kept in the master export only because it is the
input the Forest Research inventory's own pre-computed reporting fields (\texttt{Vol95},
\texttt{GYCspec95}) were built from -- it is never used as a model target or feature. An older
family built the same way from the 99th percentile (\texttt{Top\_Height99}, \texttt{Vol99},
\texttt{GYCspec99}, \texttt{elev\_percentile\_99th}) was dropped entirely from the export in an
explicit 2026-07-28 decision and is not used anywhere in this dissertation.

Beyond the target, the raw layer records a number of original stand and management fields that
were carried into the master export:

\begin{itemize}
\setlength\itemsep{0pt}
\item \textbf{Planting year} (\texttt{plyr}) and the derived \textbf{Age} (survey year minus
  planting year), used both as a modelling input and as one of the cleaning-stage filters above.
\item \textbf{Yield class} (\texttt{yldc}), a Forest Research inventory field (12 distinct values
  in the four-survey cohort, ranging 2--24) kept in the master export for audit and stratification,
  but excluded from every candidate feature set used for modelling: it is itself derived from
  forest growth information, so including it alongside a height-based target or a
  growth-curve-residual target raises a leakage/circularity concern.
\item \textbf{Windthrow hazard class} (\texttt{whcl}), an ordinal Forest Research inventory field
  (documented range 0--6; the current cohorts observe only 6 of the 7 possible levels, with class
  1 absent) recording an existing management assessment of wind-damage susceptibility, not a
  direct wind measurement.
\item \textbf{Canopy cover} (\texttt{CanopyCover}), an ALS-derived fraction (0--1) of canopy cover
  recorded per plot per survey. \fc{An internal project note
  (\texttt{documentation/data\_exploration\_gpkg/...Cheat Sheet.md}) explicitly flags
  \texttt{CanopyCover} -- along with LAI and gap fraction -- as a downstream consequence of the
  growth process itself, cautioning against its use as a predictor of top height or growth-curve
  residuals. \texttt{CanopyCover} is nonetheless currently included as one of four baseline
  features in every environmental feature set used across all three research questions
  (\texttt{documentation/methodlogy\_env\_setpick.md}, \S1.6/\S2.11). No record of a deliberate
  resolution of this tension was found in any file inspected for this chapter -- this should be
  either resolved with an explicit, documented decision, or the caution should be revisited, before
  the Methodology chapter's feature-set description is finalised.}
\item \textbf{Thinning history} (\texttt{Thin}, \texttt{last\_thinn}, \texttt{next\_thin\_}), from
  which \texttt{time\_since\_thinning} and related recency flags are derived
  (\app{exact derivation formulas and the \texttt{thinning\_status} bucketing rule}).
\item \textbf{Leaf area index} (\texttt{LAI}), kept in the master export but not used as a feature
  anywhere in this project's modelling code. \fc{LAI shows a clear processing discontinuity in
  2021 specifically -- median 3.09 and maximum 23.0 that year, against a median of 0.4--0.9 and a
  hard ceiling of 5.5 in every other survey year (verified directly on the raw layer, grouped by
  survey year). This looks like a different LAI-processing algorithm was used for the 2021
  acquisition. Since LAI is not used as a model input, this does not affect any result in this
  dissertation, but it should be stated as a caveat on the column itself, not silently omitted from
  the data dictionary.}
\item A small set of pre-computed Forest Research reporting/audit fields
  (\texttt{Vol95}, \texttt{GYCspec95}, \texttt{area}, \texttt{p1}--\texttt{p5}, \texttt{g1},
  \texttt{g3}) -- species-specific growth-model lookup coefficients and the volume/yield-class
  figures they produce, useful to foresters reading results but never used as model
  predictors, since several are direct functions of height and age
  \ct{Forest Research, "Use of Airborne Laser Scanning (ALS) for Forest Inventory" -- source of the
  p1--p5 coefficients, full reference not yet confirmed}. \app{Exact
  \texttt{Vol95}/\texttt{GYCspec95}/\texttt{Vol\_RM95} formulas.}
\end{itemize}

% FIGURE PLACEHOLDER (target and age coverage)
% Caption: "Joint distribution of stand age and ALS-derived top height across both cohorts and all
% six survey years, showing the range and density of the underlying data without presenting any
% fitted model or result."
% Content: hexbin (not scatter, to avoid overplotting at 287,064/83,382 rows) of Age vs.
% elev_percentile_95th, faceted or coloured by survey year, one panel per cohort.
\begin{figure}[!htbp]
\centering
\fbox{\parbox{0.9\textwidth}{\centering\vspace{4cm}FIGURE PLACEHOLDER --- target and age coverage,
see comment above source for exact content\vspace{4cm}}}
\caption{Age and top-height coverage across cohorts and survey years.}
\label{fig:data-target-age}
\end{figure}

\section{Derived terrain, wind, climate and soil variables}
\label{sec:data-environmental}

In addition to the fields recorded directly in the Forest Research layer, every plot in the
71,766-plot environmental table was matched to a set of external terrain, wind, climate and soil
variables, grouped below by the same category scheme used throughout the modelling pipeline
(\texttt{CATEGORY\_GROUPS} in \texttt{models/xgb\_environmental}). Full per-variable definitions,
extraction methods, and coverage statistics are recorded in
\texttt{documentation/methodlogy\_env\_setpick.md} (verified 2026-08-10 against
\texttt{data/processed/environmental/plot\_environmental\_features.parquet}); Table
\ref{tab:data-env-sources} summarises the source, resolution and temporal character of each family,
and \app{the complete per-variable valid/missing/distinct/min/max table}.

\begin{table}[!htbp]
\centering
\scriptsize
\setlength{\tabcolsep}{3pt}
\renewcommand{\arraystretch}{0.95}
\begin{tabular}{p{1.6cm}p{2.3cm}p{2.6cm}p{2.4cm}p{2.7cm}p{2.4cm}}
\hline
Data family & Source & Resolution & Temporal character & Main purpose & Plots per grid cell (avg.) \\
\hline
Forest plots & Forest Research & $\sim$20\,m $\times$ 20\,m polygons & Six survey years, 2002--2023 & Response variable, stand fields & -- (native unit) \\
Terrain & OS Terrain 50 DTM \ct{OS Terrain 50, Ordnance Survey -- exact edition/vintage not found in code, see \S\ref{sec:data-source}} & 50\,m grid & Static, single current DTM & Elevation, slope, aspect, exposure & $\sim$4.3 \\
Soil (CEH categorical) & CEH natural capital rasters \ct{CEH natural capital rasters, UKCEH} & 50\,m grid & Static & Pedotope, drainage, texture class & $\sim$4.3 \\
Soil pH & SoilGrids v2.0, ISRIC \ct{SoilGrids v2.0, ISRIC} & 250\,m grid & Static, topsoil (0--5\,cm) & Soil acidity screening & $\sim$53 \\
Wind climatology & Global Wind Atlas \ct{Global Wind Atlas} & National raster grid & Long-term climatology, static & General wind exposure & -- (raster distinct-value count varies by height/variable; 1,816--1,824 distinct values across 71,766 plots, see \texttt{methodlogy\_env\_setpick.md} \S1.2) \\
Observed wind & MIDAS Open station records \ct{MIDAS Open, Met Office} & 3 station point locations & Survey-interval summaries, 2002--2023 & Temporal wind/storm conditions & -- (point data, not gridded) \\
Climate (long-term) & CHELSA v2.1 BIOCLIM+ \ct{CHELSA v2.1} & 1\,km grid & Static, 1981--2010 climatology & GDD, precipitation & $\sim$482 \\
Climate (survey-matched) & HadUK-Grid, CEDA \ct{HadUK-Grid, CEDA/Met Office} & 1\,km grid & Time-varying, cohort-year-averaged & Temperature, ground frost & $\sim$482 \\
Spatial position & OS Open Roads/Rivers; GPKG polygons \ct{OS Open Roads, OS Open Rivers, Ordnance Survey} & Vector distance calc. & Static & Edge/context effects & -- (distance-based, not gridded) \\
\hline
\end{tabular}
\caption{External data sources matched to each plot. "Plots per grid cell" is computed directly
from plot coordinates at each resolution against the full 71,766-plot table (verified 2026-08-12);
it is an average over occupied cells, not a claim that plots are evenly distributed within a cell.
An ERA5-Land temperature extraction was also investigated (8 distinct values across all 71,766
plots, coarser than HadUK-Grid's 149) and was not retained in any final feature set.}
\label{tab:data-env-sources}
\end{table}

\textbf{Terrain.} Elevation, slope, aspect-derived northness/eastness, profile and plan curvature,
local relief, a solar-radiation index, and a frost-hollow flag are all computed from the same
50\,m OS Terrain 50 tiles. Aspect itself (a 0--360\textdegree{} compass bearing) is not used
directly as an input, since a circular bearing has a discontinuity at 0\textdegree/360\textdegree{}
that would otherwise look like a large numeric jump between two physically adjacent directions; it
is instead represented as its two continuous trigonometric components,
$\mathrm{northness}=\cos\theta$ and $\mathrm{eastness}=\sin\theta$ (\app{full derivation and the
frost-hollow flag's exact threshold rule, TPI below its own 15th percentile combined with negative
plan curvature}). Terrain values are static per plot and cannot, on their own, explain a change in
conditions between two surveys of the same plot.

\textbf{Wind.} Three distinct sources describe related but different aspects of wind exposure and
are not treated as interchangeable: the Global Wind Atlas gives a long-term, effectively static
climatology; TOPEX (computed from the same 50\,m terrain tiles, horizon angles summed over eight
compass directions) describes topographic exposure or shelter, not measured wind; and MIDAS Open
station records give temporally varying observations, summarised over the intervals between LiDAR
surveys from the three nearest stations with complete 2002--2023 metadata coverage (Glasgow
Bishopton, Strathallan Airfield, and Salsburgh -- none located at Aberfoyle itself, so station
values are a proxy for, not a measurement of, on-site conditions).

\textbf{Climate.} Two different temporal treatments exist side by side. CHELSA provides a single
static 1981--2010 climatology applied to every plot regardless of that plot's actual survey year
(2008--2023) -- a real, disclosed temporal mismatch. HadUK-Grid's temperature and ground-frost
variables are instead averaged over each cohort's own real survey years, which is closer to
survey-year-matched but still a multi-year mean, not a single matched year. At 1\,km resolution,
HadUK-Grid's temperature variable takes only 149 distinct values across the 71,766 plots -- on
average, about 482 plots share one climate-grid value, which is worth stating plainly rather than
only as a distinct-value count: a 1\,km climate cell covers roughly the same forest area as
several hundred of this study's 400\,m\textsuperscript{2} plots.

\textbf{Soil and spatial position.} SoilGrids provides topsoil pH at 250\,m resolution
(99.4\% plot coverage); three CEH categorical layers (pedotope, subsurface drainage, textural
composition) are recorded at 50\,m resolution, each with a small number of classes (3--7) shared
across many plots, and share the same 39-plot coverage gap as the CEH topographic wetness index --
one shared raster-boundary gap, not four independent failures. Distances to the nearest road,
watercourse, and forest/compartment/sub-compartment/block boundary are computed directly from the
GeoPackage's own polygon geometry and OS Open Roads/Rivers vector data; these describe spatial
context and edge position, not a direct measure of a growth process.

% FIGURE PLACEHOLDER (resolution schematic)
% Caption: "Environmental data resolution relative to plot scale. A single 400 m^2 plot sits
% inside one 50 m terrain cell alongside ~4 neighbouring plots on average, one 250 m soil cell
% alongside ~53 plots, and one 1 km climate cell alongside ~482 plots -- resolutions are not
% interchangeable and should not be read as though every plot has an independently measured
% environmental value."
% Content: nested-grid schematic, three or four concentric/overlaid squares at true relative scale
% (50 m, 250 m, 1 km) with a small dot cluster inside each representing the "plots per cell"
% averages from Table data-env-sources; a compact table variant is an acceptable fallback if the
% schematic proves hard to render clearly.
\begin{figure}[!htbp]
\centering
\fbox{\parbox{0.9\textwidth}{\centering\vspace{4cm}FIGURE PLACEHOLDER --- resolution vs.\ plot
scale schematic, see comment above source for exact content\vspace{4cm}}}
\caption{Environmental data resolution compared with plot scale.}
\label{fig:data-resolution}
\end{figure}

\section{Strengths, limitations, and what these data can support}
\label{sec:data-strengths-limitations}

The dataset's central strengths are its scale and its repeated structure: 71,766 plots in the main
cohort, drawn from a much larger 230,359-plot raw layer, with up to six ALS surveys per plot over
21 years, all concentrated on one commercially important species (Sitka spruce) with an explicit
inventory-based species filter. The six-survey cohort is a clean, fully nested subset of the
four-survey cohort rather than a partially overlapping second sample, which is what makes it a
genuine (if smaller, 48-compartment) sensitivity check on the same underlying plots. The explicit
compartment/sub-compartment/block polygons -- once the \texttt{(cpmt, scpt)} correction in
\S\ref{sec:data-hierarchy} is applied -- give a real spatial unit for compartment-level
cross-validation, and the combination of stand, management, terrain, wind, climate and soil data
matched to every plot supports a genuinely multi-source attribution analysis rather than a
height-only dataset.

Several properties limit what can be concluded from these data. Each plot contributes only four or
six observations, at unequal intervals (2--9 years), which constrains how finely a within-plot
growth trajectory can be estimated and hides whatever happens within an interval. Requiring
complete coverage across every survey year of a chosen cohort may favour plots that remained
identifiable and undisturbed throughout, rather than a uniformly random sample of the forest.
ALS-derived top height is not identical to biological growth: canopy structure, measurement
differences between acquisitions, and Forest Research's own processing choices (the confirmed 2021
LAI discontinuity is one visible example, though LAI itself is not used as a feature) can all
contribute to an apparent change between two surveys. Most of the external environmental variables
are static or long-term climatologies rather than survey-year-matched observations, and several
are recorded at a resolution coarser than the 400\,m\textsuperscript{2} plot scale -- at 1\,km, a
single climate value can describe several hundred plots at once -- so these variables are better
suited to describing broad spatial conditions than to explaining a specific change between two
surveys of the same plot. Station-based wind observations come from three stations near, but not
at, Aberfoyle. Finally, all of this is observational data: no environmental variable was
experimentally manipulated, so any association identified downstream describes a correlation, not
a causal effect.

Chapter~\ref{ch:methodology} describes how the modelling design addresses several of these
constraints directly -- for example, through compartment-level spatial-block cross-validation and
an explicit feature-screening pipeline -- but no such mechanism is assumed or explained here.
